\clearpage
\phantomsection
\section*{Declaration of Generative AI Usage}
\addcontentsline{toc}{section}{Declaration of Generative AI Usage}

I explicitly declare that during this Bachelor's Thesis, I have used generative artificial intelligence tools. AI has been part of my way of working on this project, always as an auxiliary tool and under my responsibility, but in a much deeper way than copying text into a chat and accepting an answer, as is the case for the vast majority of classmates I know and the common usage observed. This is why I wish to expand on this section.

My relationship with these tools began long before this work. In 2021 and 2022, I was already following the first advances in generative text and image models with interest. I remember generating images in Google Colab, seeing results that were still imperfect but good enough to sense that this was the beginning of something much larger and that it would have a breakthrough impact. I was fascinated to see dunes, something resembling an abstract car, or an otherwise strange scene, but generated by a computer. I also came across text models such as GPT-2, which, viewed from today's perspective, was limited and failed often, but already hinted at the significant developments to follow.

Then came more accessible text and image tools, such as ChatGPT in late 2022. I experienced it as just another development, having been involved in this field for some time, but it was not just another step; it marked the first major turning point. From then on, the way many people would write, search, code, learn, and organize work changed. A similar trend occurred with images, with noticeable improvements appearing every few months, to the point where an image created with AI, requested with proper criteria, cannot easily be distinguished from one produced by other means. This is also why I consider it important to have digital identification mechanisms, such as Google's SynthID, also implemented by other companies like OpenAI, to maintain the traceability of generated content and separate real content from that generated by AI on the internet.

In 2023, 2024, and 2025, I continued to watch generative AI transition from a curiosity into a feasible computing tool. I used ChatGPT, Google Bard, and later Gemini, Mixtral, and other emerging models, including some Chinese models that began to demonstrate that serious capabilities could be achieved with optimization. Even so, for quite a while I still viewed it as just another tool within my general interest in technology, similar to how I might be interested in new processors, graphics cards, computers, or any other technical advance.

The shift in perspective came at the end of the 2024--2025 academic year, following a talk at the School of Engineering that made me stop seeing it merely as something I followed for leisure. I understood then that I needed to truly train myself, because it was not a passing fad, but a competency that would become part of engineering work. In June 2025, I decided to learn how to use these tools with sound criteria and also to learn about them. While starting to conceptualize this Bachelor's Thesis, I took introductory courses on artificial intelligence and networks, several of them from Google, which were basic but very useful for establishing fundamentals. This training influenced the topic of this thesis and the concept of imitation control.

At that time, programming tools using AI were still in a very early stage. Claude Code had appeared in February 2025 alongside Claude 3.7 Sonnet, but it was still a niche tool, used by few and closer to completing or assisting with code than driving a full workflow. Models in the Claude Sonnet 3 family, and later Sonnet 3.5 and 3.7, were starting to become interesting for programming, but they did not yet have the capability to become a complete workspace companion. Codex also began to gain ground in 2025, first as a code agent and later with improvements for agentic programming and real-world tasks. Meanwhile, I was learning how to speak to these tools, how to prompt them, how to correct them, and how to translate my ideas into verifiable work.

I also began to learn about the entire ecosystem that was forming around agents. \textit{Agent skills} emerged, which are like small, reusable prompts or instructions to reproduce workflows. I began to understand MCPs (\textit{Model Context Protocol}), which allow connecting an AI with external tools, documentation, searches, repositories, or local systems. More integrations with editors also appeared, along with \textit{hooks} systems and ways to place boundaries on agent behavior. With all these tools, you can provide context, rules, and control gates to an AI acting freely within a folder on your computer.

During the 2025--2026 academic year, I used every assignment and lab work as an opportunity to learn how to work better with AI. I improved the way I wrote instructions, explained intent, requested reviews, and distinguished when a response was useful versus when it merely sounded convincing. I also saw the capabilities of these tools improve significantly in web design, SVG, diagrams, and figures. This was highly relevant for the thesis, as all figures are LaTeX TikZ or Python code, which can be modified, compiled, or executed and reviewed.

The second major turning point for me arrived between October and December 2025. Advanced usage was no longer about chatbots, but agents capable of planning, editing files, running commands, testing, failing, correcting, and trying again. Google Antigravity appeared in November 2025 as a development environment designed from the ground up to work with agents; OpenAI continued to push Codex as an agentic workflow tool; Claude Code and the Opus models became clear benchmarks in coding; and tools were appearing to connect agents more deeply with computer programs. In 2026, tools from the Grok Build ecosystem were also added, which have also been used in this thesis.

With all that background, I approached the development of this thesis convinced that I was going to use these tools. Not as a way to avoid work, but because it was the way of working to achieve results of much higher quality and scope. Just as it would not make sense to write a thesis in the 2010s by looking up everything in a paper encyclopedia and avoiding the internet, it did not seem reasonable to ignore tools that allow programming, documenting, reviewing, and experimenting at a much faster pace. The key was to build the project in such a way that the AI had context, documentation, rules, and constant traceability.

That is why this work was developed on my computer, using editors and tools prepared for agents. I used Google Antigravity as a code editor with integrated Gemini models, OpenAI Codex and Grok Build, NotebookLM to organize and consult documentation, standard chatbots as search tools on steroids, and the agent and code ecosystem tools mentioned above. I also used MCPs to extend the AI's capabilities, the most relevant being Context7 to have up-to-date documentation on code libraries.

What I want to make clear is the workflow: it is not 'asking a chatbot a question' and pasting it, but rather 'giving an agent a goal, context, rules, tests, and a way to verify if what it did makes sense.' When the AI only receives a snippet of code copied into a chat, it does not understand the project well, does not know past decisions, and cannot maintain traceability. In this thesis, I worked differently, with a repository set up so that the AI had context from day one and could act as my work companion.

Additionally, I used speech-recognition tools to verbalize ideas and convert them into text, as the most natural way to explain an intention is not to write it from scratch, but rather to speak it first and organize it later.

In this work, AI has been used primarily to translate intent into code and documentation. This includes Python code, configuration files, Markdown documentation, tests, commands, helper scripts, LaTeX, TikZ, diagrams, and figures. The code is not just Python; there is code in a figure made with TikZ, in a LaTeX template, a YAML file, or a JSON telemetry log. In all these cases, the AI served to accelerate the execution of decisions that I made, reviewed, and corrected.

The methodology followed was not asking 'Hello, write my thesis.' The repository was designed from the very first document to work with the AI as a companion, with specifications, plans, reviews, up-to-date documentation, requirements traceability, tests, and decision logs. The technique I followed is known as \textit{Spec-Driven Development}: before making an implementation, a specification of what is to be achieved is defined, discussed, edge cases are reviewed, and the initial intent is converted into a set of concrete requirements.

The Q\&A part with the model was central. The process started from an idea of mine, extensive but incomplete. The agent asked me questions, I answered, assumptions were corrected, my intentions were refined, thereby squaring the intent into a clear specification and plan. This recorded which idea was to be executed, what the limits were, and what decisions were made. After implementation, I reviewed the result, compared outputs, requested additional reviews, and corrected what did not fit. Thanks to the AI, we iterated faster, but the decisions regarding the scope of the work, physical extension and criteria, interpretation of results, and general direction were entirely mine.

In the public repository, all the supporting documentation generated to understand what was being done can be observed. There are plans, specifications, reviews, simulator documentation, traceability tables, requirements, maintenance instructions, working principles, and AGENTS.md files that explain to agents how they should behave within the project. People always say that AI must be given context, but I never read about anyone forcing the AI to maintain full experimental traceability.

The first iteration of the simulator was the clearest example of this methodology. I did not start with 'build a simulator for the thesis,' but with a specific intent: how I wanted the modules to be organized, what options the program should have, what type of scenarios I wanted, what data should be generated, and how I wanted to structure the execution. The AI helped convert that intent, expressed in natural language, into Python code and configuration files such as JSON, YAML, or TOML.

After that first version, the generated output was not simply accepted. There was a lot of \textit{refinement} to adjust cases, change parameters, detect bugs, improve scenarios, review outputs, run tests, rewrite parts, and verify again. The agent was left to work on the computer: editing files, launching commands, running tests, reading errors, and proposing the next fix. The AI was not answering in a chat; it was controlling the development environment to meet a verifiable objective.

Updated documentation of the simulator was created regarding architecture, usage, scenarios, telemetry, validation, requirements, testing, and the status of each part. This documentation served me, anyone who might consult the project on GitHub, and also the AI itself, because it provided context about the project. If there was a question about a module, a scenario, a metric, or a design decision, the answer depended on the documentation previously written in the repository. Thus, progress was made in an orderly manner, through multiple iterations among several agents, differentiating between implementation and review, which would have been impossible without that context.

The resulting code is structured in a modular fashion and with software best practices, but without obfuscating the important part of the work: physics, aerospace, and control. Functions related to dynamics, reference frames, integration, control, or network training have been documented with clear names, comments, and explanations.

The AI was also used to create auxiliary tools that are not the core of the simulator but facilitate working with it: scripts, commands, analysis utilities, results generation, documentation review, and support for reproducibility.

I have also used AI in the thesis manuscript to support essentially programming-like tasks. The writing—that is, the intent of what is to be said and defended—is entirely mine. But LaTeX is code, and it is written in the same editor as Python. That is why the AI worked with the preambles, table and figure parameters, references, layout, and general figure handling.

The latter are a particularly important case. None of them are images generated by image-generative AI; they are diagrams made with TikZ (LaTeX) and Python, which is code that is modified and executed like any other. I used Codex and Gemini to generate, adjust, and refine these figures, \textbf{always} starting from my own intent: what the diagram should show, what elements it should contain, and a proposed layout and color palette. All figures have a Markdown sheet detailing their intent, conventions, and traceability.

Without these tools, I would not have been able to make figures with this level of quality and flexibility. The AI understands TikZ and LaTeX well because, at their core, they are code. Thus, I have been able to include custom-made figures in the work instead of generic images. In my opinion, this shows the real potential: allowing anyone, including all researchers and scientists, to use them for artifacts that they would otherwise not make. The expert defines the what and the why; the AI fills in the how.

Another relevant aspect of the work has been verbalization. I have often used speech as the primary interface to explain ideas, as it is the most natural way to communicate intent. This is not exactly generative AI, but rather speech recognition and workflow integration, yet it was important for this thesis. I used my own local tools, such as Voice Stall, to transcribe voice to text and express long ideas before converting them into plans or instructions for agents.

In conclusion, yes, I have used generative AI for programming, documentation, information retrieval, linguistic coherence, figure generation and tuning, code review, repository analysis, script writing, LaTeX work, refinement of specifications with plans, and execution of agentic workflows. I assume full responsibility for the originality, veracity, rigor, and authorship of this document. In my view, this usage has been part of my learning as a future Aerospace Engineer and represents a professional competency that will be in high demand. Thanks to using it this way during the academic year, I can explain which tools I use, the procedures, the limits, the review process, and how AI can be used without sacrificing human originality. I can present this repository and not just claim that I used AI, but show how, thereby building trust with others to share knowledge.

Just as it would make no sense in 2015 to write a paper ignoring the internet and searching in paper encyclopedias, it does not seem reasonable to ignore tools that are already a real part of engineering work in multiple companies, most prominently in the United States. The difference lies in the usage: a substitute or a capability multiplier. In this thesis, I have used it as the latter. For me, this is the path forward: not to ban or demonize generative AI, but to use it responsibly to increase the complexity of what can be achieved, program better with less experience, and produce high-quality material.
